\documentclass[12pt,a4paper]{report}
\usepackage[utf8]{inputenc}
\usepackage[T1]{fontenc}
\usepackage[french]{babel}
\usepackage{geometry}
\geometry{top=2.5cm,bottom=2.5cm,left=2.5cm,right=2.5cm}
\usepackage{graphicx}
\usepackage{xcolor}
\usepackage{tikz}
\usetikzlibrary{shapes.geometric,arrows.meta,positioning,calc,shadows.blur,fit,backgrounds}
\usepackage{tcolorbox}
\tcbuselibrary{skins,breakable}
\usepackage{fontawesome5}
\usepackage{enumitem}
\usepackage{hyperref}
\hypersetup{colorlinks=true,linkcolor=indigo,urlcolor=indigo,pdfauthor={MonEcole},pdftitle={Manuel — Module Recouvrement}}
\usepackage{fancyhdr}
\usepackage{tabularx}
\usepackage{booktabs}
\usepackage{float}
\usepackage{caption}

% ====== COULEURS ======
\definecolor{indigo}{HTML}{4338CA}
\definecolor{violet}{HTML}{7C3AED}
\definecolor{amber}{HTML}{D97706}
\definecolor{emerald}{HTML}{059669}
\definecolor{rose}{HTML}{E11D48}
\definecolor{slate}{HTML}{334155}
\definecolor{slatelight}{HTML}{94A3B8}
\definecolor{teal}{HTML}{0D9488}
\definecolor{sky}{HTML}{0EA5E9}
\definecolor{coverbg}{HTML}{1E1B4B}
\definecolor{coveraccent}{HTML}{059669}
\definecolor{lightbg}{HTML}{F8FAFC}
\definecolor{warningbg}{HTML}{FEF3C7}
\definecolor{successbg}{HTML}{ECFDF5}
\definecolor{infobg}{HTML}{EFF6FF}
\definecolor{dangerbg}{HTML}{FEF2F2}
\definecolor{purplebg}{HTML}{FAF5FF}
\definecolor{tealbg}{HTML}{F0FDFA}

% ====== TCOLORBOX ======
\newtcolorbox{astuce}[1][]{enhanced,colback=successbg,colframe=emerald,fonttitle=\bfseries\small,title={\faLightbulb\hspace{6pt}Astuce},left=8pt,right=8pt,top=6pt,bottom=6pt,arc=4pt,boxrule=1pt,breakable,#1}
\newtcolorbox{attention}[1][]{enhanced,colback=warningbg,colframe=amber,fonttitle=\bfseries\small,title={\faExclamationTriangle\hspace{6pt}Attention},left=8pt,right=8pt,top=6pt,bottom=6pt,arc=4pt,boxrule=1pt,breakable,#1}
\newtcolorbox{info}[1][]{enhanced,colback=infobg,colframe=indigo,fonttitle=\bfseries\small,title={\faInfoCircle\hspace{6pt}Information},left=8pt,right=8pt,top=6pt,bottom=6pt,arc=4pt,boxrule=1pt,breakable,#1}
\newtcolorbox{danger}[1][]{enhanced,colback=dangerbg,colframe=rose,fonttitle=\bfseries\small,title={\faTimesCircle\hspace{6pt}Important},left=8pt,right=8pt,top=6pt,bottom=6pt,arc=4pt,boxrule=1pt,breakable,#1}
\newtcolorbox{etape}[1][]{enhanced,colback=white,colframe=indigo!60,fonttitle=\bfseries\small,left=10pt,right=10pt,top=8pt,bottom=8pt,arc=6pt,boxrule=1.5pt,breakable,shadow={2pt}{-2pt}{0pt}{black!12},#1}
\newtcolorbox{prerequis}[1][]{enhanced,colback=purplebg,colframe=violet,fonttitle=\bfseries\small,title={\faClipboard\hspace{6pt}Pr\'erequis},left=8pt,right=8pt,top=6pt,bottom=6pt,arc=4pt,boxrule=1pt,breakable,#1}
\newtcolorbox{resultat}[1][]{enhanced,colback=tealbg,colframe=teal,fonttitle=\bfseries\small,title={\faCheck\hspace{6pt}R\'esultat attendu},left=8pt,right=8pt,top=6pt,bottom=6pt,arc=4pt,boxrule=1pt,breakable,#1}

% ====== HEADERS ======
\pagestyle{fancy}\fancyhf{}
\fancyhead[L]{\small\color{slatelight}\leftmark}
\fancyhead[R]{\small\color{emerald}\textbf{MonEcole}}
\fancyfoot[C]{\small\color{slatelight}\thepage}
\renewcommand{\headrulewidth}{0.5pt}
\renewcommand{\footrulewidth}{0.3pt}

% ====== COMMANDES ======
\newcommand{\bouton}[1]{{\colorbox{emerald}{\small\sffamily\bfseries\color{white}\,#1\,}}}
\newcommand{\onglet}[1]{{\colorbox{indigo}{\small\sffamily\bfseries\color{white}\,#1\,}}}
\newcommand{\champ}[1]{\texttt{\color{violet}#1}}
\newcommand{\btnviolet}[1]{{\colorbox{violet}{\small\sffamily\bfseries\color{white}\,#1\,}}}
\newcommand{\btnamber}[1]{{\colorbox{amber}{\small\sffamily\bfseries\color{white}\,#1\,}}}
\newcommand{\btnrose}[1]{{\colorbox{rose}{\small\sffamily\bfseries\color{white}\,#1\,}}}

\begin{document}

% ===================================================================
%                         PAGE DE COUVERTURE
% ===================================================================
\begin{titlepage}
\begin{tikzpicture}[remember picture,overlay]
  \fill[coverbg] (current page.south west) rectangle (current page.north east);
  \fill[coveraccent,opacity=0.12] (current page.north east) ++(-4,1) circle (8cm);
  \fill[teal,opacity=0.08] (current page.south west) ++(3,-2) circle (10cm);
  \fill[coveraccent,opacity=0.06] (current page.center) ++(6,-4) circle (5cm);
  \fill[amber,opacity=0.05] (current page.north west) ++(2,-8) circle (4cm);
  \draw[coveraccent,line width=2pt,opacity=0.4] ([yshift=-6cm]current page.north west) ++(2.5,0) -- ++(16,0);
  \node[fill=coveraccent,rounded corners=14pt,minimum size=2cm,inner sep=0pt] at ([yshift=-3.5cm]current page.north) {\color{white}\fontsize{28pt}{28pt}\selectfont\faMoneyBillWave};
  \node[anchor=north,text width=14cm,align=center] at ([yshift=-6.5cm]current page.north) {
    \color{white}\fontsize{32pt}{36pt}\selectfont\textbf{Manuel d'Utilisation}\\[12pt]
    \fontsize{18pt}{22pt}\selectfont\color{coveraccent!70!white}Module \textbf{Recouvrement}
  };
  \node[anchor=north,text width=14cm,align=center] at ([yshift=-10cm]current page.north) {
    \color{white!60}\fontsize{11pt}{15pt}\selectfont
    Guide complet pour la gestion financi\`ere scolaire~:\\
    configuration des variables de paiement, suivi des paiements,\\
    p\'enalit\'es de retard, dates butoires et op\'erations de caisse.
  };
  \draw[coveraccent,line width=1.5pt] ([yshift=-12.5cm]current page.north) ++(-3,0) -- ++(6,0);
  \node[anchor=north,text width=14cm,align=center] at ([yshift=-13.5cm]current page.north) {
    \color{white!50}\small
    \faBuilding\hspace{4pt}Plateforme MonEcole --- Architecture Hub-and-Spoke\\[6pt]
    \faCalendar\hspace{4pt}Mai 2026 --- Version 1.0\\[6pt]
    \faBook\hspace{4pt}Documentation Utilisateur --- 4 Sections D\'etaill\'ees
  };
  \node[anchor=south,text width=14cm,align=center] at ([yshift=1.5cm]current page.south) {
    \color{white!30}\footnotesize \faCopyright\hspace{3pt} Nexora Tech --- Tous droits r\'eserv\'es
  };
\end{tikzpicture}
\end{titlepage}

\tableofcontents
\newpage

% ===================================================================
\chapter{Introduction au Module Recouvrement}
% ===================================================================

\section{Pr\'esentation g\'en\'erale}

Le module \textbf{Recouvrement} est le centre de gestion financi\`ere de la plateforme \textbf{MonEcole}. Il couvre l'ensemble du cycle de vie des paiements scolaires~: de la d\'efinition des frais jusqu'\`a la g\'en\'eration des re\c{c}us, en passant par le suivi en temps r\'eel, les r\'eductions, les p\'enalit\'es de retard et la gestion des op\'erations de caisse.

Ce module est accessible depuis la barre de navigation lat\'erale (sidebar) en cliquant sur \textbf{Recouvrement}. Il se divise en \textbf{4 sections} principales~:

\begin{center}
\begin{tikzpicture}[
  sect/.style={rounded corners=10pt, minimum height=1.1cm, minimum width=3.5cm, font=\small\bfseries, inner sep=8pt, text=white, drop shadow={shadow xshift=1pt, shadow yshift=-1pt, opacity=0.2}},
  desc/.style={text width=3.3cm, align=center, font=\tiny\color{slate}}
]
  \node[sect, fill=violet] (s1) at (0,0) {\faChartPie\ Dashboard};
  \node[sect, fill=emerald, right=6pt of s1] (s2) {\faMoneyBillWave\ Paiements};
  \node[sect, fill=amber, right=6pt of s2] (s3) {\faCogs\ Configuration};
  \node[sect, fill=rose, right=6pt of s3] (s4) {\faCashRegister\ Caisse};

  \node[desc, below=8pt of s1] {Vue d'ensemble financi\`ere avec graphiques et statistiques};
  \node[desc, below=8pt of s2] {Enregistrement et suivi des paiements par \'el\`eve};
  \node[desc, below=8pt of s3] {Cat\'egories, variables, prix, p\'enalit\'es, dates butoires, banques};
  \node[desc, below=8pt of s4] {Entr\'ees et sorties de caisse hors paiements scolaires};
\end{tikzpicture}
\end{center}

\section{Flux de travail recommand\'e}

Pour un \'etablissement nouvellement configur\'e, voici l'ordre logique des op\'erations~:

\begin{center}
\begin{tikzpicture}[
  step/.style={draw=slate!40, fill=white, rounded corners=8pt, minimum width=2.2cm, minimum height=1.5cm, align=center, font=\scriptsize\bfseries, text=slate, drop shadow={shadow xshift=1pt, shadow yshift=-1pt, opacity=0.12}},
  arrow/.style={-{Stealth[length=6pt,width=5pt]}, thick, color=slatelight}
]
  \node[step] (s1) at (0,0) {\faUniversity\\[2pt]Cr\'eer les\\banques};
  \node[step, right=0.4cm of s1] (s2) {\faFolder\\[2pt]Cr\'eer les\\cat\'egories};
  \node[step, right=0.4cm of s2] (s3) {\faTags\\[2pt]D\'efinir les\\variables};
  \node[step, right=0.4cm of s3] (s4) {\faMoneyCheckAlt\\[2pt]Fixer les\\prix};
  \node[step, right=0.4cm of s4] (s5) {\faMoneyBillWave\\[2pt]Enregistrer\\paiements};
  \node[step, right=0.4cm of s5] (s6) {\faChartPie\\[2pt]Suivre via\\dashboard};

  \draw[arrow] (s1) -- (s2);
  \draw[arrow] (s2) -- (s3);
  \draw[arrow] (s3) -- (s4);
  \draw[arrow] (s4) -- (s5);
  \draw[arrow] (s5) -- (s6);
\end{tikzpicture}
\end{center}

\begin{danger}
Il est \textbf{imp\'eratif} de suivre cet ordre. Vous ne pouvez pas enregistrer de paiements si aucun prix n'a \'et\'e d\'efini, et vous ne pouvez pas d\'efinir de prix sans avoir cr\'e\'e les variables et les banques/comptes.
\end{danger}

\section{Concepts cl\'es}

\begin{description}[leftmargin=1.5cm, style=nextline, font=\bfseries\color{indigo}]
  \item[Cat\'egorie de variable] Regroupement logique des frais (ex.\ : Minerval, Fournitures, Activit\'es).
  \item[Variable] Type de frais sp\'ecifique (ex.\ : Frais scolaire T1, Uniforme, Transport).
  \item[Prix] Montant \`a payer pour une variable dans une classe donn\'ee et une ann\'ee donn\'ee.
  \item[R\'eduction] Pourcentage de remise accord\'e \`a un \'el\`eve sp\'ecifique sur une variable.
  \item[Date butoire] Date limite de paiement pour une variable dans une classe.
  \item[D\'erogation] Prolongation individuelle de la date butoire pour un \'el\`eve.
  \item[P\'enalit\'e] Montant forfaitaire ou pourcentage appliqu\'e en cas de retard de paiement.
  \item[Banque / Compte] Institution et num\'ero de compte bancaire pour les encaissements.
  \item[Op\'eration de caisse] Entr\'ee ou sortie financi\`ere hors paiements scolaires.
\end{description}

\newpage
% ===================================================================
\chapter{Section 1 --- Dashboard Financier}
% ===================================================================

\section{Objectif de cette section}

Le \textbf{Dashboard} offre une vue d'ensemble de la situation financi\`ere de l'\'etablissement. Il affiche des indicateurs cl\'es et des graphiques permettant de suivre la progression des recouvrements.

\section{Les indicateurs cl\'es}

Six cartes statistiques sont affich\'ees en haut de la page~:

\begin{center}
\begin{tabularx}{\textwidth}{llX}
\toprule
\textbf{Indicateur} & \textbf{Couleur} & \textbf{Description} \\
\midrule
Total Pay\'e & Vert \'emeraude & Somme de tous les paiements valid\'es \\
Total Attendu & Ambre & Somme des prix configur\'es $\times$ nombre d'\'el\`eves inscrits \\
Reste \`a Payer & Rose & Diff\'erence entre le total attendu et le total pay\'e \\
Transactions & Violet & Nombre total de paiements enregistr\'es \\
En Dette & Bleu & Nombre d'\'el\`eves ayant un solde impay\'e \\
Rejet\'es & Teal & Nombre de paiements rejet\'es par l'administration \\
\bottomrule
\end{tabularx}
\end{center}

\begin{etape}[title=Proc\'edure : Consulter le dashboard]
\begin{enumerate}[leftmargin=1.5cm,label=\textbf{\arabic*.},itemsep=10pt]
  \item \textbf{Acc\'edez au Recouvrement} en cliquant sur \textbf{Recouvrement} dans la barre lat\'erale.
  \item \textbf{Cliquez sur Dashboard} (ic\^one \faChartPie) si ce n'est pas la section active.
  \item \textbf{S\'electionnez l'ann\'ee} dans le menu d\'eroulant \champ{Ann\'ee} en haut \`a droite du graphique. Les donn\'ees se rechargent automatiquement.
  \item \textbf{Consultez la barre de progression} qui indique le pourcentage de recouvrement global.
  \item \textbf{Consultez le r\'esum\'e} dans le panneau lat\'eral droit (transactions, \'el\`eves en dette, rejet\'es).
\end{enumerate}
\end{etape}

\begin{astuce}
Les cartes statistiques sont cliquables. Cliquez sur une carte pour voir les d\'etails correspondants dans le tableau en dessous du graphique.
\end{astuce}

\newpage
% ===================================================================
\chapter{Section 2 --- Gestion des Paiements}
% ===================================================================

\section{Objectif de cette section}

La section \textbf{Paiements} est le c\oe ur op\'erationnel du module. Elle permet d'enregistrer les paiements des \'el\`eves, de les valider, de les rejeter et de g\'en\'erer des fiches de paiement en PDF.

\section{Comprendre l'interface}

L'interface se compose de~:

\begin{enumerate}[leftmargin=1.2cm]
  \item \textbf{Barre de s\'electeurs} --- Trois menus d\'eroulants en cascade~:
  \begin{itemize}
    \item \champ{Ann\'ee} --- S\'electionnez l'ann\'ee scolaire.
    \item \champ{Classe} --- Se charge apr\`es s\'election de l'ann\'ee.
    \item \champ{\'El\`eve} --- Se charge apr\`es s\'election de la classe.
  \end{itemize}
  \item \textbf{Bouton} \bouton{Payer} --- Appara\^it apr\`es s\'election d'un \'el\`eve.
  \item \textbf{Cartes de situation} --- Variables avec barre de progression (pay\'e / reste).
  \item \textbf{Onglets de paiements} --- En attente, Valid\'es, Rejet\'es.
\end{enumerate}

\section{Pour enregistrer un nouveau paiement}

\begin{prerequis}
\begin{itemize}
  \item Les \textbf{variables} et \textbf{prix} doivent \^etre configur\'es (section Configuration).
  \item Des \textbf{banques et comptes} doivent exister.
  \item L'\textbf{\'el\`eve} doit \^etre inscrit dans la classe pour l'ann\'ee s\'electionn\'ee.
\end{itemize}
\end{prerequis}

\begin{etape}[title=Proc\'edure d\'etaill\'ee : Enregistrer un paiement]
\begin{enumerate}[leftmargin=1.5cm,label=\textbf{\arabic*.},itemsep=12pt]
  \item \textbf{S\'electionnez l'ann\'ee} dans le menu \champ{Ann\'ee}.
  \item \textbf{S\'electionnez la classe} dans le menu \champ{Classe}. La liste des \'el\`eves se charge.
  \item \textbf{S\'electionnez l'\'el\`eve} dans le menu \champ{\'El\`eve}. Les cartes de situation (variables avec reste \`a payer) s'affichent et le bouton \bouton{Payer} appara\^it.
  \item \textbf{Cliquez sur} \bouton{Payer}. Le formulaire de paiement s'affiche.
  \item \textbf{S\'electionnez la variable} dans le menu \champ{Variable}. Seules les variables avec un reste \`a payer sont list\'ees, avec le montant restant affich\'e.
  \item \textbf{Saisissez le montant} dans le champ \champ{Montant}. Le syst\`eme v\'erifie que le montant ne d\'epasse pas le reste \`a payer.
  \item \textbf{S\'electionnez la date} de paiement dans \champ{Date paiement}. Les dates futures sont refus\'ees.
  \item \textbf{S\'electionnez le compte} bancaire dans \champ{Compte}.
  \item \textbf{Joignez le bordereau} (optionnel) --- Cliquez sur \champ{Bordereau} pour joindre une photo ou PDF du re\c{c}u.
  \item \textbf{Cliquez sur} \bouton{Enregistrer}. Un toast vert confirme la r\'eussite.
\end{enumerate}
\end{etape}

\begin{resultat}
Le paiement appara\^it dans l'onglet \textbf{En attente}. Les cartes de situation de l'\'el\`eve se mettent \`a jour automatiquement avec le nouveau montant pay\'e et le reste \`a payer recalcul\'e.
\end{resultat}

\begin{attention}
Un paiement enregistr\'e a le statut \textit{En attente} par d\'efaut. Il doit \^etre \textbf{valid\'e} par un administrateur pour \^etre pris en compte dans les statistiques et les fiches officielles.
\end{attention}

\section{Pour valider ou rejeter un paiement}

\begin{etape}[title=Proc\'edure : Valider / Rejeter]
\begin{enumerate}[leftmargin=1.5cm,label=\textbf{\arabic*.},itemsep=10pt]
  \item \textbf{Acc\'edez \`a l'onglet} \onglet{En attente} dans la section Paiements.
  \item \textbf{Rep\'erez le paiement} dans le tableau (nom \'el\`eve, variable, montant, date).
  \item \textbf{Pour valider}~: cliquez sur \textcolor{emerald}{\faCheck} (bouton vert). Le paiement passe dans l'onglet Valid\'es.
  \item \textbf{Pour rejeter}~: cliquez sur \textcolor{rose}{\faTimes} (bouton rouge). Le paiement passe dans l'onglet Rejet\'es.
\end{enumerate}
\end{etape}

\section{Pour g\'en\'erer une fiche de paiement PDF}

\begin{etape}[title=Proc\'edure : Exporter les fiches PDF]
\begin{enumerate}[leftmargin=1.5cm,label=\textbf{\arabic*.},itemsep=10pt]
  \item \textbf{S\'electionnez la classe} dans les s\'electeurs.
  \item \textbf{Pour la fiche classe}~: cliquez sur \btnrose{Fiche Classe} (bouton PDF rouge). Un PDF r\'ecapitulatif de tous les paiements de la classe s'ouvre dans un nouvel onglet.
  \item \textbf{Pour la fiche \'el\`eve}~: s\'electionnez l'\'el\`eve puis cliquez sur \btnviolet{Fiche \'El\`eve}. Un PDF individuel s'ouvre.
  \item \textbf{Pour une facture}~: dans l'onglet Valid\'es, cliquez sur \textcolor{rose}{\faFilePdf} \`a c\^ot\'e d'un paiement pour g\'en\'erer la facture unitaire.
\end{enumerate}
\end{etape}

\newpage
% ===================================================================
\chapter{Section 3 --- Configuration Financi\`ere}
% ===================================================================

\section{Objectif de cette section}

La section \textbf{Configuration} permet de param\'etrer l'ensemble des \'el\'ements financiers de l'\'etablissement. Elle est organis\'ee en \textbf{6 onglets}~:

\begin{center}
\begin{tikzpicture}[
  btn/.style={rounded corners=4pt, minimum height=0.6cm, font=\tiny\bfseries, inner sep=4pt, text=white, fill=#1}
]
  \node[btn=indigo] (b1) at (0,0) {\faFolder\ Cat\'egories};
  \node[btn=violet, right=3pt of b1] (b2) {\faTags\ Variables};
  \node[btn=teal, right=3pt of b2] (b3) {\faMoneyCheckAlt\ Prix};
  \node[btn=amber, right=3pt of b3] (b4) {\faExclamationTriangle\ P\'enalit\'es};
  \node[btn=rose, right=3pt of b4] (b5) {\faCalendarTimes\ Dates Butoires};
  \node[btn=slate, right=3pt of b5] (b6) {\faUniversity\ Banques};
\end{tikzpicture}
\end{center}

% manuel_recouvrement_suite.tex — inclus par manuel_recouvrement.tex

% -------- Onglet Cat\'egories --------
\section{Onglet 1 --- Cat\'egories de Variables}

Les cat\'egories permettent de regrouper les frais par famille logique. Exemples~: \textit{Minerval}, \textit{Fournitures}, \textit{Activit\'es parascolaires}.

\begin{etape}[title=Proc\'edure : Cr\'eer une cat\'egorie]
\begin{enumerate}[leftmargin=1.5cm,label=\textbf{\arabic*.},itemsep=10pt]
  \item \textbf{Cliquez sur l'onglet} \onglet{Cat\'egories} dans la barre de configuration.
  \item \textbf{Remplissez le champ} \champ{Nom} dans le formulaire en haut (ex.\ : \textit{Minerval}).
  \item \textbf{Cliquez sur} \bouton{Ajouter}. La cat\'egorie appara\^it dans le tableau ci-dessous.
\end{enumerate}
\end{etape}

\begin{resultat}
La cat\'egorie est imm\'ediatement disponible dans le menu d\'eroulant \champ{Cat\'egorie} lors de la cr\'eation de variables (onglet suivant).
\end{resultat}

% -------- Onglet Variables --------
\section{Onglet 2 --- Variables de Paiement}

Les variables repr\'esentent les \'el\'ements concrets \`a payer. Chaque variable est rattach\'ee \`a une cat\'egorie.

\begin{etape}[title=Proc\'edure : Cr\'eer une variable]
\begin{enumerate}[leftmargin=1.5cm,label=\textbf{\arabic*.},itemsep=10pt]
  \item \textbf{Cliquez sur l'onglet} \onglet{Variables}.
  \item \textbf{S\'electionnez la cat\'egorie} dans le menu \champ{Cat\'egorie}.
  \item \textbf{Saisissez le nom} dans \champ{Nom variable} (ex.\ : \textit{Frais scolaire T1}).
  \item \textbf{Cliquez sur} \bouton{Ajouter}. La variable appara\^it dans le tableau.
\end{enumerate}
\end{etape}

\begin{info}
Chaque variable dispose d'un bouton \textbf{Obligatoire} (Oui/Non) qui permet de marquer les frais obligatoires. Cliquez sur le bouton vert/orange dans la colonne \textbf{Obligatoire} pour basculer l'\'etat.
\end{info}

% -------- Onglet Prix --------
\section{Onglet 3 --- Prix par Variable et Classe}

C'est ici que vous d\'efinissez le montant exact que chaque classe doit payer pour chaque variable. C'est le c\oe ur de la configuration financi\`ere.

\begin{etape}[title=Proc\'edure d\'etaill\'ee : D\'efinir les prix]
\begin{enumerate}[leftmargin=1.5cm,label=\textbf{\arabic*.},itemsep=12pt]
  \item \textbf{S\'electionnez l'ann\'ee} dans le menu \champ{Ann\'ee}.
  \item \textbf{S\'electionnez la classe} dans le menu \champ{Classe}. Le tableau des variables se charge.
  \item \textbf{Pour chaque variable}, saisissez le prix dans le champ num\'erique de la colonne \champ{Prix (Fbu)}.
  \item \textbf{Cliquez sur} \bouton{Enregistrer} sur la ligne correspondante.
  \item Le badge passe de \textcolor{amber}{\textit{Non d\'efini}} \`a \textcolor{emerald}{\textit{D\'efini}}.
\end{enumerate}
\end{etape}

\begin{center}
\begin{tikzpicture}
  \fill[white, rounded corners=6pt, drop shadow={shadow xshift=1pt, shadow yshift=-1pt, opacity=0.1}] (0,0) rectangle (13,3.2);
  \draw[indigo!20, rounded corners=6pt, line width=1pt] (0,0) rectangle (13,3.2);
  % Header
  \fill[slate!8, rounded corners=4pt] (0.2,2.4) rectangle (12.8,2.9);
  \node[font=\tiny\bfseries\color{slate}] at (1,2.65) {N\textdegree};
  \node[font=\tiny\bfseries\color{slate}] at (3.5,2.65) {Variable};
  \node[font=\tiny\bfseries\color{slate}] at (6.5,2.65) {Cat\'egorie};
  \node[font=\tiny\bfseries\color{slate}] at (9.5,2.65) {Prix (Fbu)};
  \node[font=\tiny\bfseries\color{slate}] at (11.8,2.65) {Action};
  % Row 1
  \node[font=\tiny\color{slate}] at (1,2) {1};
  \node[font=\tiny\bfseries\color{slate}] at (3.5,2) {Frais scolaire T1};
  \node[fill=emerald!12, rounded corners=2pt, inner sep=2pt, font=\tiny\color{emerald}] at (6,2) {\faCheck\ D\'efini};
  \node[font=\tiny\color{slate}] at (7.2,2) {Minerval};
  \node[draw=slate!20, rounded corners=3pt, minimum width=1.5cm, font=\tiny\color{slate}, inner sep=2pt] at (9.5,2) {150 000};
  \node[fill=emerald!80, rounded corners=3pt, inner sep=2pt, font=\tiny\bfseries\color{white}] at (11.8,2) {\faSave\ Enregistrer};
  % Row 2
  \node[font=\tiny\color{slate}] at (1,1.3) {2};
  \node[font=\tiny\bfseries\color{slate}] at (3.5,1.3) {Uniforme};
  \node[fill=amber!12, rounded corners=2pt, inner sep=2pt, font=\tiny\color{amber}] at (6,1.3) {\faExclamationCircle\ Non d\'efini};
  \node[font=\tiny\color{slate}] at (7.2,1.3) {Fournitures};
  \node[draw=slate!20, rounded corners=3pt, minimum width=1.5cm, font=\tiny\color{slatelight}, inner sep=2pt] at (9.5,1.3) {0};
  \node[fill=emerald!80, rounded corners=3pt, inner sep=2pt, font=\tiny\bfseries\color{white}] at (11.8,1.3) {\faSave\ Enregistrer};
  % Summary
  \node[font=\tiny\color{slate}, anchor=west] at (0.4,0.5) {\textcolor{emerald}{\faCheckCircle}\ \textbf{1} prix d\'efini \quad \textcolor{amber}{\faExclamationCircle}\ \textbf{1} en attente};
\end{tikzpicture}
\captionof{figure}{Exemple du tableau de configuration des prix}
\end{center}

\begin{attention}
Les prix sont sp\'ecifiques \`a chaque couple \textbf{ann\'ee + classe}. Si vous changez d'ann\'ee ou de classe, vous devez red\'efinir les prix. Cela permet une flexibilit\'e maximale (prix diff\'erents selon les niveaux).
\end{attention}

% -------- Onglet P\'enalit\'es --------
\section{Onglet 4 --- R\`egles de P\'enalit\'e}

Les p\'enalit\'es permettent de configurer des frais suppl\'ementaires automatiques en cas de retard de paiement. Deux types sont disponibles~:

\begin{center}
\begin{tabularx}{\textwidth}{lX}
\toprule
\textbf{Type} & \textbf{Description} \\
\midrule
Forfait & Un montant fixe ajout\'e au solde de l'\'el\`eve (ex.\ : 5\,000 Fbu de p\'enalit\'e). \\
Pourcentage & Un pourcentage du montant restant d\^u (ex.\ : 10\% du reste \`a payer). \\
\bottomrule
\end{tabularx}
\end{center}

\begin{etape}[title=Proc\'edure : Cr\'eer une r\`egle de p\'enalit\'e]
\begin{enumerate}[leftmargin=1.5cm,label=\textbf{\arabic*.},itemsep=10pt]
  \item \textbf{Cliquez sur l'onglet} \onglet{P\'enalit\'es}.
  \item \textbf{S\'electionnez l'ann\'ee} dans \champ{Ann\'ee}.
  \item \textbf{S\'electionnez la variable} concern\'ee dans \champ{Variable} (ou laissez sur \textit{Toutes} pour une p\'enalit\'e globale).
  \item \textbf{Choisissez le type}~: \textit{Forfait} ou \textit{Pourcentage (\%)}.
  \item \textbf{Saisissez la valeur} (montant en Fbu pour un forfait, pourcentage pour un \%).
  \item \textbf{Saisissez le plafond} (optionnel) --- montant maximal que la p\'enalit\'e ne peut pas d\'epasser.
  \item \textbf{Cliquez sur} \bouton{Ajouter}. La r\`egle appara\^it dans le tableau ci-dessous.
\end{enumerate}
\end{etape}

% -------- Onglet Dates Butoires --------
\section{Onglet 5 --- Dates Butoires}

Les dates butoires d\'efinissent les dates limites de paiement par variable et par classe.

\begin{etape}[title=Proc\'edure : D\'efinir une date butoire]
\begin{enumerate}[leftmargin=1.5cm,label=\textbf{\arabic*.},itemsep=10pt]
  \item \textbf{Cliquez sur l'onglet} \onglet{Dates Butoires}.
  \item \textbf{S\'electionnez l'ann\'ee} et la \textbf{classe}.
  \item \textbf{S\'electionnez la variable} concern\'ee.
  \item \textbf{Saisissez la date limite} dans \champ{Date limite}.
  \item \textbf{Cliquez sur} \bouton{Enregistrer}.
\end{enumerate}
\end{etape}

\begin{info}
Les dates d\'epass\'ees apparaissent en \textcolor{rose}{\textbf{rouge}} dans le tableau, et les dates futures en \textcolor{emerald}{\textbf{vert}}. Cela permet de rep\'erer rapidement les \'ech\'eances d\'epass\'ees.
\end{info}

% -------- Onglet Banques --------
\section{Onglet 6 --- Banques et Comptes}

Cet onglet est divis\'e en deux panneaux c\^ote \`a c\^ote~:

\begin{enumerate}[leftmargin=1.2cm]
  \item \textbf{Panneau gauche} --- \textbf{Banques}~: institutions bancaires (ex.\ : BANCOBU, KCB).
  \item \textbf{Panneau droit} --- \textbf{Comptes bancaires}~: num\'eros de compte li\'es \`a une banque.
\end{enumerate}

\begin{etape}[title=Proc\'edure : Ajouter une banque et un compte]
\begin{enumerate}[leftmargin=1.5cm,label=\textbf{\arabic*.},itemsep=10pt]
  \item \textbf{Cliquez sur l'onglet} \onglet{Banques \& Comptes}.
  \item \textbf{Dans le panneau Banques}~: saisissez le \champ{Nom} et le \champ{Sigle}, puis cliquez sur \bouton{+}.
  \item \textbf{Dans le panneau Comptes}~: s\'electionnez la \champ{Banque} dans le menu d\'eroulant, saisissez le \champ{N\textdegree\ compte}, puis cliquez sur \bouton{+}.
\end{enumerate}
\end{etape}

\begin{resultat}
Les comptes bancaires sont imm\'ediatement disponibles dans le menu d\'eroulant \champ{Compte} du formulaire de paiement (section Paiements).
\end{resultat}

\newpage
% ===================================================================
\chapter{Section 4 --- Entr\'ees et Sorties de Caisse}
% ===================================================================

\section{Objectif de cette section}

La section \textbf{Caisse} permet de suivre les flux financiers \textbf{hors paiements scolaires}~: achats de fournitures, salaires, dons re\c{c}us, frais d'entretien, etc. Elle est compl\`etement ind\'ependante du cycle de paiement des \'el\`eves.

\section{Les indicateurs de caisse}

Quatre cartes statistiques en haut~:

\begin{center}
\begin{tabularx}{\textwidth}{llX}
\toprule
\textbf{Indicateur} & \textbf{Couleur} & \textbf{Description} \\
\midrule
Total Op\'erations & Violet & Nombre total d'op\'erations enregistr\'ees \\
Total Entr\'ees & Vert & Somme de toutes les entr\'ees (recettes) \\
Total Sorties & Rose & Somme de toutes les sorties (d\'epenses) \\
Solde & Bleu & Diff\'erence Entr\'ees $-$ Sorties \\
\bottomrule
\end{tabularx}
\end{center}

\section{Pour cr\'eer une cat\'egorie d'op\'eration}

\begin{etape}[title=Proc\'edure : Cr\'eer une cat\'egorie d'op\'eration]
\begin{enumerate}[leftmargin=1.5cm,label=\textbf{\arabic*.},itemsep=10pt]
  \item \textbf{Cliquez sur l'onglet} \onglet{Cat\'egories d'op\'erations} dans la section Caisse.
  \item \textbf{S\'electionnez le type}~: \textit{Entr\'ee} ou \textit{Sortie}.
  \item \textbf{Saisissez le nom} (ex.\ : \textit{Fournitures bureau}).
  \item \textbf{Ajoutez une description} (optionnel).
  \item \textbf{Cliquez sur} \bouton{Ajouter}.
\end{enumerate}
\end{etape}

\section{Pour enregistrer une op\'eration}

\begin{etape}[title=Proc\'edure d\'etaill\'ee : Enregistrer une op\'eration de caisse]
\begin{enumerate}[leftmargin=1.5cm,label=\textbf{\arabic*.},itemsep=12pt]
  \item \textbf{Acc\'edez \`a l'onglet} \onglet{Journal des op\'erations}.
  \item \textbf{Cliquez sur} \bouton{Nouvelle op\'eration}. Le formulaire s'affiche.
  \item \textbf{S\'electionnez le type}~: \textit{Entr\'ee} ou \textit{Sortie}. Les cat\'egories correspondantes se chargent.
  \item \textbf{S\'electionnez la cat\'egorie} d'op\'eration.
  \item \textbf{Saisissez le montant} (positif obligatoirement).
  \item \textbf{Saisissez la date} de l'op\'eration.
  \item \textbf{S\'electionnez le mode} de paiement~: Esp\`eces, Ch\`eque, Virement, Autre.
  \item \textbf{Compl\'etez} les champs optionnels~: Source/B\'en\'eficiaire, R\'ef\'erence, Description.
  \item \textbf{Cliquez sur} \bouton{Enregistrer}.
\end{enumerate}
\end{etape}

\section{Pour exporter le journal}

\begin{etape}[title=Proc\'edure : Exporter les op\'erations]
\begin{enumerate}[leftmargin=1.5cm,label=\textbf{\arabic*.},itemsep=8pt]
  \item \textbf{Filtrez} les op\'erations par ann\'ee, cat\'egorie ou type si souhait\'e.
  \item \textbf{Pour un PDF}~: cliquez sur \btnrose{PDF} en haut \`a droite du journal.
  \item \textbf{Pour un Excel}~: cliquez sur \bouton{Excel} en haut \`a droite. Le fichier se t\'el\'echarge directement.
\end{enumerate}
\end{etape}

\newpage
% ===================================================================
\chapter{R\'ef\'erence rapide}
% ===================================================================

\section{Navigation entre les sections}

\begin{tabularx}{\textwidth}{llX}
\toprule
\textbf{Ic\^one Sidebar} & \textbf{Section} & \textbf{Description} \\
\midrule
\faChartPie\ (violet) & Dashboard & Tableau de bord financier avec indicateurs et graphiques \\
\faMoneyBillWave\ (vert) & Paiements & Enregistrement, validation et suivi des paiements \\
\faCogs\ (ambre) & Configuration & Cat\'egories, variables, prix, p\'enalit\'es, banques \\
\faCashRegister\ (rouge) & Caisse & Entr\'ees et sorties financi\`eres hors scolaire \\
\bottomrule
\end{tabularx}

\section{Endpoints API du module Recouvrement}

{\small
\begin{tabularx}{\textwidth}{lX}
\toprule
\textbf{Endpoint} & \textbf{Description} \\
\midrule
\texttt{/api/recouvrement/dashboard-data/} & GET --- Statistiques g\'en\'erales du dashboard \\
\texttt{/api/recouvrement/classes/} & GET --- Classes actives pour une ann\'ee \\
\texttt{/api/recouvrement/eleves/} & GET --- \'El\`eves inscrits dans une classe \\
\texttt{/api/recouvrement/variables-restant/} & GET --- Variables avec reste \`a payer d'un \'el\`eve \\
\texttt{/api/recouvrement/save-paiement/} & POST --- Enregistrer un paiement \\
\texttt{/api/recouvrement/paiements-submitted/} & GET --- Paiements en attente \\
\texttt{/api/recouvrement/paiements-validated/} & GET --- Paiements valid\'es et rejet\'es \\
\texttt{/api/recouvrement/save-categorie-variable/} & POST --- Cr\'eer une cat\'egorie \\
\texttt{/api/recouvrement/save-variable/} & POST --- Cr\'eer une variable \\
\texttt{/api/recouvrement/save-variable-prix/} & POST --- D\'efinir un prix \\
\texttt{/api/recouvrement/save-penalite/} & POST --- Cr\'eer une r\`egle de p\'enalit\'e \\
\texttt{/api/recouvrement/save-date-butoire/} & POST --- D\'efinir une date butoire \\
\texttt{/api/recouvrement/save-banque/} & POST --- Cr\'eer une banque \\
\texttt{/api/recouvrement/save-compte/} & POST --- Cr\'eer un compte bancaire \\
\texttt{/api/recouvrement/operations/} & GET --- Journal des op\'erations de caisse \\
\texttt{/api/recouvrement/save-operation/} & POST --- Enregistrer une op\'eration \\
\texttt{/api/recouvrement/invoice/<id>/} & GET --- G\'en\'erer une facture PDF \\
\bottomrule
\end{tabularx}
}

\section{Glossaire}

\begin{description}[leftmargin=1.5cm, style=nextline, font=\bfseries\color{indigo}]
  \item[Bordereau] Pi\`ece justificative (photo ou PDF) jointe \`a un paiement.
  \item[Cat\'egorie] Regroupement logique de variables de paiement.
  \item[Date butoire] Date limite de paiement pour une variable donn\'ee.
  \item[D\'erogation] Prolongation individuelle d'une date butoire pour un \'el\`eve.
  \item[Fiche classe] Document PDF r\'ecapitulatif des paiements de toute une classe.
  \item[Fiche \'el\`eve] Document PDF individuel d\'etaillant les paiements d'un \'el\`eve.
  \item[Op\'eration de caisse] Mouvement financier (entr\'ee ou sortie) hors scolaire.
  \item[P\'enalit\'e] Frais suppl\'ementaire appliqu\'e en cas de retard de paiement.
  \item[R\'eduction] Remise accord\'ee \`a un \'el\`eve (en pourcentage du prix total).
  \item[Reste \`a payer] Montant total d\^u moins les paiements valid\'es.
  \item[Variable] \'El\'ement de facturation sp\'ecifique (frais scolaire, uniforme, etc.).
\end{description}

\vspace{2cm}
\begin{center}
\begin{tikzpicture}
  \fill[coverbg, rounded corners=12pt] (0,0) rectangle (14,2.5);
  \node[text=white, font=\large\bfseries] at (7,1.5) {\faCheck\hspace{8pt}Fin du Manuel --- Module Recouvrement};
  \node[text=white!50, font=\small] at (7,0.7) {Plateforme MonEcole --- Nexora Tech --- Mai 2026 --- v1.0};
\end{tikzpicture}
\end{center}


\end{document}
